\section{Optical weights and transport at a critical roof}
\label{sec:v76-optical}
The next estimate is on the actual regular word, without a lower
incidence cutoff. It improves an exponential cardinality majorant
only at the level of its optical weights; it does not estimate the
complete word sum.

For a regular word of $m$ flights let $u,v$ be oriented endpoint
arclengths, $F$ its reduced length, $c_i>0$ its incidences, and
$\ell_i\ge\tau_-$ its flight lengths. Put $t_i=\ell_i^{-1}$ and
$V_i=2/(Rc_i)$ for $1\le i<m$. The internal optical matrix is
\begin{equation}\label{eq:v76-optical-matrix}
 A_{ii}=t_{i-1}+t_i+V_i,\qquad A_{i,i+1}=A_{i+1,i}=-t_i.
\end{equation}
Signs of internal tangents can be chosen successively to give this
form. The full internal action Hessian is $D_cAD_c$, with
$D_c=\operatorname{diag}(c_1,\ldots,c_{m-1})$.

\begin{lemma}[Transfer denominator and monotonicity]
\label{lem:v76-transfer}
Define $y_0=0$, $f_0=1$, $y_1=\ell_0$, and, for $1\le i<m$,
\begin{equation}\label{eq:v76-transfer-recursion}
 f_i=f_{i-1}+V_i y_i,\qquad y_{i+1}=y_i+\ell_i f_i.
\end{equation}
Then
\begin{equation}\label{eq:v76-denominator}
 y_m=\left(\prod_{i=0}^{m-1}\ell_i\right)\det A,
 \qquad |F_{uv}|=\frac{c_0c_m}{y_m}.
\end{equation}
The empty determinant is one. The positive polynomial $y_m$ is
nondecreasing in every flight length and every internal potential.
\end{lemma}
\begin{proof}
Eliminating $f_i$ gives
\[
 y_{i+1}=(1+\ell_i/\ell_{i-1}+\ell_iV_i)y_i
                              -(\ell_i/\ell_{i-1})y_{i-1}.
\]
This is exactly the continuant recursion for
$(\prod_{j=0}^{i}\ell_j)\det A_{[1,i]}$, with the same initial
values. It proves the first identity, including $m=1$.
The mixed endpoint derivative is the corner Schur complement of
the contact Hessian. Its numerator is the product of neighboring
mixed entries, and the internal factors $c_i^2$ cancel from its
determinant. Thus
$|F_{uv}|=c_0c_m\prod t_i/\det A$, proving the second identity.
The positive recursion \eqref{eq:v76-transfer-recursion}, rather
than its subtractive second-order form, proves the monotonicity
and the positive-polynomial assertion by induction.
\end{proof}

Put
\begin{equation}\label{eq:v76-optical-rate}
 \eta=\operatorname{arcosh}(53/47),\qquad
 \rho=e^{-\eta}=\frac{53-10\sqrt6}{47}<1,
 \qquad Q(c)=\frac{53c}{6+47c}.
\end{equation}
The guard parameter $q_0=47/53$ has not been changed to $\rho$.

\begin{theorem}[Explicit inverse twist with simultaneous incidence gains]
\label{thm:v76-optical-weight}
Let $B_m=\partial_{p_0}\alpha_m$. For every subset
$S\subset\{1,\ldots,m-1\}$,
\begin{equation}\label{eq:v76-optical-bound}
 \frac1{|B_m|}=R|F_{uv}|
 \le\frac{R c_0c_m}{\tau_-}
       \frac{\sinh\eta}{\sinh(m\eta)}\prod_{i\in S}Q(c_i).
\end{equation}
All constants are independent of the physical word, its length,
positive incidences, and positive clearances. In particular each
bounded-source Gauss summand of the preceding geometric theorem
is bounded by a fixed constant times
\[
 Wc_0c_m\frac{\sinh\eta}{\sinh(m\eta)}\prod_{i\in S}Q(c_i),
\]
with its exact event indicator still present.
\end{theorem}
\begin{proof}
First take $S$ empty. Monotonicity in Lemma~\ref{lem:v76-transfer}
allows all $\ell_i$ to be replaced by $\tau_-$ and all $V_i$ by
$2/R_+$. The resulting recurrence is
$y_{i+1}=2(1+\tau_-/R_+)y_i-y_{i-1}$, with $y_0=0$ and
$y_1=\tau_-$. Since $1+\tau_-/R_+=53/47$, its solution is
\[
 y_m^{\min}=\tau_-\frac{\sinh(m\eta)}{\sinh\eta}.
\]
The inverse-coordinate Jacobian and $u=R\alpha_0$, $v=R\alpha_m$
give $|B_m|^{-1}=R|F_{uv}|$. This proves the empty-set estimate.

For general $S$, let $A^{(S)}$ replace $c_i$ by one in the selected
diagonal potentials only. Restore the selected incidences one at
a time. Before restoring $c_i$, its scalar Schur complement $s_i$
is positive and at most $2/\tau_-+2/R$. The determinant ratio for
that step is
\[
 \frac{s_i}{s_i+(2/R)(c_i^{-1}-1)}
 \le\frac{(1+R/\tau_-)c_i}{1+(R/\tau_-)c_i}
 \le Q(c_i).
\]
Each intermediate matrix is positive definite, so multiplication
gives $\det A^{(S)}/\det A\le\prod_{i\in S}Q(c_i)$.
The empty-set transfer bound also applies to $A^{(S)}$, since all
its potentials remain at least $2/R_+$. This proves
\eqref{eq:v76-optical-bound}. Finally the preceding Gauss formula
bounds its summand by
$C_GW\one_\Pi\one_a/|B_m|$, where
$C_G=2R_+/(\pi c_*R_-^2)$. Substitution gives the last assertion.
No comparison matrix has been assigned a physical orbit or source.
\end{proof}

\subsection{Interior critical points and the direction of the jump}
A regular interior critical point of the roof has $p_0=p_m=0$ by
the first variation. Let $K=D^2_{(\alpha_0,\alpha_m)}F$ there.
At that point
\begin{equation}\label{eq:v76-critical-hessian}
                         K\ge R I_2.
\end{equation}
Indeed the affine support supplied by its own physical trajectory
restricts to the two endpoint circles as
$F_*+R(1-\cos(\alpha_0-\alpha_{0,*}))
     +R(1-\cos(\alpha_m-\alpha_{m,*}))$.
It touches the physical roof from below. Differentiating twice at
the contact proves \eqref{eq:v76-critical-hessian}. Internal strict
positivity gives smoothness on a regular physical neighborhood.

\begin{proposition}[An intrinsic positive critical charge]
\label{prop:v76-critical-charge}
At a regular interior critical point $x_*$ whose neighborhood stays
in the exact event, the local density for a continuous locally
Lipschitz nonnegative weight $W$ has right-minus-left jump
\begin{equation}\label{eq:v76-critical-charge}
 J_*^W=\frac{W(x_*)}{2c_*|B_m(x_*)|\sqrt{\det K}}\ge0.
\end{equation}
For every selected internal set $S$ it satisfies
\begin{equation}\label{eq:v76-critical-charge-bound}
 J_*^W\le\frac{W(x_*)}{2c_*\tau_-}
              \frac{\sinh\eta}{\sinh(m\eta)}
                        \prod_{i\in S}Q(c_i(x_*)).
\end{equation}
These are local charges. Other physical or decision-boundary
contributions at the same roof are not asserted to have this sign.
\end{proposition}
\begin{proof}
In endpoint angles the source density is
$a_\alpha=W/(4\pi c_*|B_m|)$. Smooth Morse coordinates $z\in\R^2$
put $F=F_*+|z|^2/2$; their inverse Jacobian at zero is
$(\det K)^{-1/2}$. Coarea in polar coordinates shows that the
right limiting density is $2\pi a_\alpha(x_*)/\sqrt{\det K}$,
while the local source below $F_*$ is zero. A cutoff equal to one
near $x_*$ does not change this jump. This proves
\eqref{eq:v76-critical-charge}. Since $\sqrt{\det K}\ge R$,
$ c_0=c_m=1$, and \eqref{eq:v76-optical-bound} applies, the stated
bound follows. The assertion concerns the local nonnegative source,
not the Jordan decomposition of unrelated terms at the same roof.
\end{proof}

\begin{proposition}[Morse transport does not charge a positive jump]
\label{prop:v76-morse-transport}
Use the preceding Morse coordinates on a disk $|z|<r_*$, and write
the full weighted physical area there as $a(z)\dd z$, where $a$
is nonnegative and Lipschitz with constant $L_a$. The forward roof
flow
\begin{equation}\label{eq:v76-morse-flow}
 z\longmapsto z_h=\sqrt{1+2h/|z|^2}\,z
\end{equation}
has Euclidean Jacobian one and increases $F$ by exactly $h$.
If $0<t-F_*<r_*^2/8$ and $0<h<r_*^2/8$, it stays in this
physical disk and its unmatched initial mass satisfies
\begin{equation}\label{eq:v76-morse-defect}
                  R_h^+(t)\le2\pi L_a\sqrt{2h}.
\end{equation}
Here the level source is the local disk source, so no exterior
trajectory is included. In particular the coefficient $J_*^W$
itself is not an exit or distortion charge in this direction.
\end{proposition}
\begin{proof}
Put $r=|z|$ and $r_h=\sqrt{r^2+2h}$. The radial and tangential
derivative factors are $r/r_h$ and $r_h/r$, respectively, and their
product is one. The full circle at the starting roof and its image
lie in the disk under the stated restrictions. Polar coarea for
$|z|^2/2$ is $\dd\theta$. The unmatched mass is consequently
\[
 \int_0^{2\pi}[a(r e_\theta)-a(r_h e_\theta)]_+\dd\theta
 \le2\pi L_a(r_h-r)\le2\pi L_a\sqrt{2h}.
\]
The constant part of $a$ makes zero contribution, including its
positive jump at $F_*$. Push this vector field back to the original
collision manifold. It satisfies $XL=1$ away from the critical point,
so Theorem~\ref{thm:v76-subcoupling} applies. A smooth partition of
unity glues this field to the normalized gradient outside a smaller
disk, since a convex combination of fields satisfying $XL=1$ still
satisfies that identity. There are finitely many interior critical
points at fixed count, and disjoint neighborhoods can be used.
Neither their radii nor the sum of their Lipschitz constants is
asserted uniform in the collision count.
\end{proof}
